\chapter{Structure of Complete Market}

\section{General Discrete-Time Markets}

Up until this point, we have meticulously analyzed a highly specialized discrete market containing exactly $d=1$ risky asset and $1$ riskless asset, constrained to exhibit strictly binomial price transitions. A natural theoretical question arises: Are there more generalized discrete-time frameworks that naturally yield complete markets? 

To rigorously investigate this structural property, consider a generalized $T$-period market modeled upon a probability space $(\Omega, \mathcal{F}, \mathbb{P})$. The market contains one riskless asset serving cleanly as the numeraire, and an arbitrary multidimensional vector of $d$ risky assets modeled collectively as a $d$-dimensional stochastic process $S = (S^1, \dots, S^d)$. 

The explicit transition dynamics of these multiple assets follow:
\begin{equation}
    \Delta S_t := S_t - S_{t-1} \in \mathbb{R}^d.
\end{equation}
The underlying multidimensional market is \textit{complete} if and only if every bounded $\mathcal{F}_T$-measurable random variable can be replicated via a self-financing strategy across all $d$ risky assets and the singular bank account unit. Driven by the Second Fundamental Theorem of Asset Pricing, this holds if and exclusively if the equivalent martingale measure $\mathbb{Q}$ acts uniquely.

\section{The \texorpdfstring{$d+1$}{d+1} Branching Theorem}

The geometric constraints governing market completeness lock branching probability distributions toward low dimensional geometries matching the underlying dimension of liquid financial instruments.

\begin{proposition}[Geometric Branching Structure of Complete Markets]
    Assume that there exists a \textbf{unique} martingale measure (the market is arbitrage-free and complete). 
    
    Then, there exists an adapted process consisting of exactly $d+1$ vectors $\{\delta_{t}^i\}_{i=1}^{d+1}$ formally taking values in $\mathbb{R}^d$, tightly enforcing computationally that:
    \begin{equation} \label{eq:branching_condition}
        \sum_{i=1}^{d+1} \mathbb{P}\left(\Delta S_t = \delta_{t-1}^i \;\middle|\; \mathcal{F}_{t-1}\right) = 1 \quad \mathbb{P}\text{-a.s.} \;\; \forall t \in \{1, \dots, T\}.
    \end{equation}
    Moreover, for all $t$, the family of state vectors $\{\delta_t^1, \dots, \delta_t^{d+1}\}$ is strictly \textbf{affinely independent} $\mathbb{P}$-a.s. Mathematically, this enforces that the derivative vectors:
    \begin{equation}
        \{ \delta_t^2 - \delta_t^1, \dots, \delta_t^{d+1} - \delta_t^1 \}
    \end{equation}
    are \textbf{linearly independent} vectors strongly spanning $\mathbb{R}^d$ almost surely.
\end{proposition}

\begin{proof}
    The geometric mappings dominating multi-asset pricing heavily rely upon balancing the equations dynamically. At any node, if the path branches into $K$ disjoint possibilities, we generate $K$ target portfolio values to match. Our strategy permits choosing explicit allocations spanning exactly $d$ risky assets supplemented alongside $1$ riskless bank account. 
    Therefore, perfect replication yields a system of $K$ equations with exactly $d+1$ structural unknowns. Resolving this unique geometric parity demands $K \leq d+1$. The linear independence rigidly ensures full spatial span to prevent redundant paths, conclusively enforcing uniquely solvable algorithmic replication.
\end{proof}

\section{The Singularity of the Binomial Model}

This abstract spatial theorem immediately generates powerful analytical consequences regarding discrete-time modeling constraints.

Substituting identical values for $d=1$ smoothly generates exactly $1+1=2$ explicitly valid branching paths. Thus:
\begin{enumerate}
    \item A non-trivial model prevents pure market stagnation, generating bounded price shifts: $\mathbb{P}[S_t = S_{t-1}] < 1$.
    \item The functionally necessary matching of exactly $2$ independent branches fiercely limits complete 1-dimensional modeling flawlessly to the Cox-Ross-Rubinstein \textbf{binomial model}.
\end{enumerate}

\begin{corollary}[Incompleteness of Complex Trees]
    Consequently, excluding isolated tree arrays precisely configured to span $d+1$ branches for exactly $d$ independent liquid assets, generalized discrete time structural models naturally and decisively exhibit pure \textbf{incompleteness}.
\end{corollary}

In explicit financial vocabulary, heavily attempting to utilize trinomial (3-branch) trees to map completely hedgeable options operating upon merely a single ($d=1$) underlying asset predictably fails. The excess unknown states relative to the explicit trading tools natively embeds unhedgable risk structurally operating fundamentally beyond replication entirely.
